\section{应用实验与结果解释}
实验单位是具有明确工作负载、软件修订与网络配置的一次调用或应用会话。对于路径可观测的状态，分别报告提交、准入、选定、完成、拒绝、超时和取消。成功率需要明确分母及截止时间。延迟分布应说明包含哪些请求，避免较快的已完成请求掩盖更大的失败比例。重复实验或独立随机种子支持变异估计；单次探索运行应如实标记。

UAV 会话记录命令完成延迟、视频启动延迟、有效帧交付、录制获取完成及操作员看到的状态年龄。多视角工作还记录可用视角数、处理完成与基于标注观测的分析指标。框架实验首先固定识别算法；声称多视角提升识别准确率需要独立应用研究。

DI 记录模型准备、prefill、首 token 时间、后续输出速率、总完成时间、设备峰值内存及 Provider 间传输字节。Cold 与 warm 分别比较，保留可运行的单 Provider 对照。如果分布式部署源于模型无法装入单设备，应直接报告容量收益，而非相对不可运行基线构造加速比。通信、序列化、密码操作和模型计算分项剖析用于定位瓶颈。

授权评价应覆盖双方服务权限，而不只是解密成功。反例包括无权限调用者、无权限 Provider、签名有效但受保护值错误的消息，以及请求关闭后的重复使用。授予实验测量已登记实体如何发现并应用新权限；撤销实验测量更新传播、状态验证、参数和 DKEY 刷新，以及应用更新后的拒绝。断连节点和旧密文分别检验，并遵守授权论证中说明的限制。

协作消融保留 request-ID 隔离和签名验证。去除指定生产者检查时，测试同一请求内的错误角色对象；去除依赖边检查时，测试来自获准参与者的未声明输入；替换测试注入重新分配后的旧工作。这能隔离扩展的作用，而不把已有调用基础解决的跨请求混淆再次声称为创新。正确性测试需要预期拒绝及禁止的 handler 未执行，而不只是日志出现错误。

框架复用可通过记录 UAV 与 DI 共享的调用、授权、对象传输和生命周期组件，以及仍需应用适配的部分来评价。代码行数不足以衡量可维护性；还需记录增加服务类型所需修改和必须提供的契约。这支撑工程理由，协议正确性与性能则由各自证据支持。

部署记录应包含依赖版本、模型摘要、拓扑、链路参数、时钟、重试策略、缓存及密钥管理配置，但不公开秘密。应用失败应与启动、配置和结果收集失败分开；环境未建立并不能证明协议无法完成请求。这些记录帮助最终论文区分受支持结论、有价值的负结果和未测试配置。
